Sources : Insee, André et Meslin (2023), Insee Analyses n° 91, données de la figure 5a (taxe foncière 2017 / valeur de marché estimée, par commune) ; DGFiP, DVF 2024 (prix médian au m² des ventes d'un seul logement) ; calculs de l'auteur sur les 4 661 communes comptant au moins 20 ventes.\par Lecture : dans les 5 \% de communes les moins chères (prix médian 925 \euro{}/m²), la taxe foncière représente en médiane 0,54 \% de la valeur des logements chaque année ; dans les 5 \% les plus chères (5 900 \euro{}/m²), 0,20 \%. Dans chaque vingtième, la moitié des communes se situe dans la bande claire. La taxe date de 2017 et les prix de 2024 : on suppose que le classement des communes par prix a peu changé entre les deux dates.
